Artificial-life systems can grow, move and preserve organized structures through local interactions. A further challenge is to make a continuing body learn from the consequences of its own actions. This requires more than an interesting trajectory: experience must change internal content, that content must contribute to subsequent behaviour, and the body must remain able to obtain and use new evidence. When material also flows, consumes energy and disappears, learning and physical persistence become coupled experimental questions.

Continuous cellular automata make these questions concrete. Lenia and its extensions produce localized and interacting structures, while parameter search discovers sensorimotor behaviours\cite{chan,expandedlenia,plantec,hamon,asal}. Neural cellular automata support growth and regeneration\cite{growingca,developmental,bricks}. Runtime plasticity, meta-learning and differentiable updates provide routes towards adaptation during operation\cite{khajehabdollahi,pdnca,duan,wang,miconi,backpropamine}. These achievements do not, by themselves, distinguish a rule selected by an external optimizer from content acquired within an individual's lifetime. Nor does a history-dependent response establish an advantage over an equally informed fixed physiological rule.

The distinction is especially important when information is supplied by the experimental interface. A controller receiving a reward gradient already knows a useful direction. A system observing completed displacements can estimate an execution map unavailable to a scalar-outcome learner. A recent-feedback rule may solve a delayed association task if no competing event intervenes. Conversely, a learner can fail because its own behaviour prevents contact with new evidence. Evaluating endpoints without specifying these interfaces confounds acquired capability with the observations and actions installed in advance.

Studies of non-neural organisms motivate treating memory as a physical process rather than assigning it from appearance alone. Habituation, environmental traces, network reorganization and retained responses can all connect past conditions to future behaviour\cite{boisseau,reid,kramar2021,vogel,boussard,reiber}. Regulatory and collective dynamics offer further forms of persistence\cite{acar,micheelsen,yang,lambert}. Transporting a state with tissue, however, does not establish that its content is learned or necessary. Regeneration adds a separate requirement: behavioural use must persist when its original material carrier is actually lost\cite{shomrat}.

Here we organize these distinctions into three experimental tests. First, we specify what the learner can observe and compare it with simple rules using the same lawful information. Second, we branch complete acquired bodies and manipulate only the proposed content, retaining physical state and randomness at the branch. Third, we measure the resulting movement, energy, injury and tissue production. Strong controls and physical accounting are essential because learning can be costly and its benefit depends on environmental timing\cite{mery2005,dunlap2009,mitchell,tagkopoulos}. Matched branches test mechanisms within histories; they do not create additional independent learning histories.

We apply this framework to three explicitly separated implementations: conservative moving bodies for action and cue learning, resource-constrained generative bodies for admission and persistence, and an externally supplied density-regulated birth/death model for behaviour during tissue production and predictive maintenance. The main results proceed from scalar action acquisition to controlled temporal credit, then to spatial consequences and conditional physical utility. Seven figures connect behavioural tests, actual fields and resource accounts. Standard learning rules serve as experimental instruments; the contribution is the resolved relationship between acquired content and physical outcomes, rather than a claim to a new general optimizer. The models jointly identify requirements for a regenerating learner, but their separate positive results are not combined into evidence that one such learner has already been achieved.
